\documentclass[12pt,a4paper]{article}

% ==========================
% == PACCHETTI PRINCIPALI ==
% ==========================
\usepackage[utf8]{inputenc}
\usepackage[T1]{fontenc}
\usepackage[italian]{babel}
\usepackage{geometry}
\usepackage{graphicx}
\usepackage{lmodern}
\usepackage{setspace}
\usepackage{fancyhdr}
\usepackage{titlesec}
\usepackage{hyperref}
\usepackage[table]{xcolor}
\usepackage{tabularx}
\usepackage{booktabs}
\usepackage{microtype}
\usepackage{enumitem}
\usepackage{multirow}
\usepackage{tcolorbox}
\usepackage{lastpage}
\usepackage{amsmath}

% ==========================
% == CONFIGURAZIONE BASE ===
% ==========================
\geometry{margin=2.5cm}
\onehalfspacing
\setcounter{tocdepth}{3}
\setcounter{secnumdepth}{3}

% === Colori ===
\definecolor{darkgray}{gray}{0.25}
\definecolor{lightgray}{gray}{0.95}
\definecolor{uniblue}{HTML}{003366}

% === Stile titoli ===
\titleformat{\section}{\normalfont\Large\bfseries\color{uniblue}}{\thesection.}{1em}{}
\titleformat{\subsection}{\normalfont\large\bfseries\color{uniblue}}{\thesubsection.}{1em}{}
\titleformat{\subsubsection}{\normalfont\normalsize\bfseries\color{uniblue}}{\thesubsubsection.}{1em}{}

% === Link nel PDF ===
\hypersetup{
	colorlinks=true,
	linkcolor=uniblue,
	urlcolor=uniblue,
	pdftitle={Documentazione - Adaptive Bitboard Engine},
	pdfauthor={Gruppo Manfredini-Romano},
}

% ==========================
% ==== HEADER & FOOTER =====
% ==========================
\fancypagestyle{standard}{
	\fancyhf{}
	\fancyfoot[L]{\small\textbf{UNISA – Corso di Fondamenti d'intelligenza Artificiale}}
	\fancyfoot[R]{\small\textbf{Pag. \thepage \quad | \quad \pageref{LastPage}}}
	
	\fancyhead[L]{\small\textbf{Adaptive Bitboard Engine | Introduzione}}
	\fancyhead[R]{\small\textit{\nouppercase{\leftmark}}}
	
	\renewcommand{\headrulewidth}{0.4pt}
	\renewcommand{\footrulewidth}{0.4pt}
}

\renewcommand{\sectionmark}[1]{\markboth{#1}{}}

% ==========================
% ====== FRONTESPIZIO ======
% ==========================
\begin{document}
	\thispagestyle{empty}
	
	% --- Intestazione istituzionale con logo UNISA e testo a sinistra ---
	\begin{minipage}{0.75\textwidth}
		{\LARGE \textbf{Università degli Studi di Salerno}}\\
		{\large Corso di Ingegneria del Software}
	\end{minipage}
	\hfill
	\begin{minipage}{0.2\textwidth}
		\raggedleft
		% Logica robusta: Se il file esiste lo mostra, altrimenti mostra un box placeholder
		\IfFileExists{images/logo_unisa.png}{%
			\includegraphics[width=3cm]{images/logo_unisa.png}%
		}{%
			\framebox{\parbox[c][3cm][c]{3cm}{\centering\small\textbf{LOGO\\UNISA}}}%
		}
	\end{minipage}
	\\[0.2cm]
	{\color{uniblue}\rule{\textwidth}{0.4pt}}
	\vspace{1.8cm}
	
	% --- Logo del progetto centrale (grande) ---
	\begin{center}
		% Logica robusta anche per il logo progetto
		\IfFileExists{images/logo_Forza_IV.png}{%
			\includegraphics[width=7cm]{images/logo_Forza_IV.png}%
		}{%
			\framebox{\parbox[c][4cm][c]{7cm}{\centering\textbf{LOGO PROGETTO}}}\\[1cm]%
		}
		
		{\Huge \textbf{Adaptive Bitboard Engine}}\\[0.5cm]
		{\Large \textbf{Documentazione di Progetto - Versione 1.0}}\\[2cm]
		
		\vspace{1cm}
		\begin{tcolorbox}[colback=lightgray!10, colframe=uniblue,
			title=Informazioni Generali, fonttitle=\bfseries]
			\renewcommand{\arraystretch}{1.3}
			\begin{tabularx}{\textwidth}{l X}
				\textbf{Corso di Laurea:} & Informatica \\
				\textbf{Università:} & Università degli Studi di Salerno (UNISA) \\
				\textbf{Docente:} & Chiar.mo Prof. PALOMBA \\
				\textbf{Data:} & 30/01/2026 \\
				\textbf{Anno Accademico:} & 2025/2026 \\
			\end{tabularx}
		\end{tcolorbox}
		
	\end{center}
	
	\clearpage
	\thispagestyle{empty}
	
	% ==========================
	% === MEMBRI DEL GRUPPO ===
	% ==========================
	
	\vspace*{1cm}
	{\Large\bfseries\color{uniblue} Membri Del Gruppo}\\[-0.2cm]
	\color{uniblue}\rule{\textwidth}{0.4pt}\\[0.6cm]
	
	\renewcommand{\arraystretch}{1.3}
	\setlength{\tabcolsep}{10pt}
	
	\begin{center}
		\begin{tcolorbox}[colback=lightgray!10, colframe=uniblue, fonttitle=\bfseries]
			\begin{tabularx}{\textwidth}{l X}
				\textbf{MANFREDINI Umberto} & Matricola 0512119797 \\
				\textbf{ROMANO Pino Fiorello} & Matricola 0512120259 \\
			\end{tabularx}
		\end{tcolorbox}
	\end{center}
	
	\vspace{1cm}
	
	% ==========================
	% === REVISION HISTORY ===
	% ==========================
	
	\vspace*{1cm}
	{\Large\bfseries\color{uniblue} Revision History}\\[-0.2cm]
	\color{uniblue}\rule{\textwidth}{0.4pt}\\[0.6cm]
	
	\renewcommand{\arraystretch}{1.3}
	
	\begin{center}
		\rowcolors{2}{lightgray!30}{white}
		\begin{tabularx}{\textwidth}{|c|c|X|c|}
			\hline
			\rowcolor{uniblue!20}
			\textbf{Data} & \textbf{Versione} & \textbf{Descrizione} & \textbf{Autore} \\
			\hline
			30/01/2026 & 1.0 & Stesura Capitolo 1 (Introduzione e Visione) & Gruppo \\
			\hline
			-- & -- & -- & -- \\
			\hline
		\end{tabularx}
	\end{center}
	
	\clearpage
	
	\newpage
	\pagestyle{standard}
	
	% ==========================
	% ======== INDICE ==========
	% ==========================
	
	\tableofcontents
	\newpage
	
	% =============================================
	% ===== CAPITOLO 1: INTRODUZIONE ==============
	% =============================================
	
	\section{Introduzione e Visione Progettuale}
	
	\subsection{Abstract}
	Il gioco del Forza 4 (\textit{Connect Four}) rappresenta, nel campo dell'Intelligenza Artificiale, un problema di ricerca spaziale a informazione perfetta risolto matematicamente. Tuttavia, le implementazioni tradizionali basate esclusivamente sulla Teoria dei Giochi classica tendono a produrre agenti ``statici'', incapaci di adattarsi allo stile specifico dell'avversario umano.
	
	Il presente progetto propone lo sviluppo di un motore decisionale ibrido, denominato \textbf{Adaptive Bitboard Engine}. Il sistema integra un algoritmo di ricerca \textit{Minimax} con \textit{Alpha-Beta Pruning} ad alte prestazioni (basato su rappresentazione \textit{Bitboard}) con un modulo di \textbf{Profilazione Comportamentale}. L'obiettivo è superare il concetto di ``mossa ottimale assoluta'' per approdare alla ``mossa ottimale relativa'', ovvero quella che massimizza la probabilità di vittoria sfruttando i bias cognitivi e le debolezze tattiche dello specifico avversario.
	
	\subsection{Limiti della Teoria dei Giochi Classica}
	Nella sua formulazione classica, un agente basato su \textit{Minimax} opera sotto l'assunto di razionalità perfetta: si presuppone che l'avversario (il giocatore \textit{Minimizing}) esegua sempre la mossa migliore possibile per contrastare l'agente (\textit{Maximizing}).
	La funzione di valore $V(s)$ per uno stato $s$ è definita ricorsivamente come:
	
	\begin{equation}
		V(s) = \max_{a \in A} \left( \min_{b \in B} f(s') \right)
	\end{equation}
	
	Dove $A$ e $B$ sono gli insiemi delle mosse possibili per i due giocatori. Sebbene questo approccio garantisca che l'IA non perda (o perda il più tardi possibile) contro un avversario perfetto, presenta evidenti limiti nell'interazione uomo-macchina:
	
	\begin{itemize}
		\item \textbf{Rigidità Strategica:} L'agente ignora le opportunità derivanti da errori non forzati dell'avversario se questi non rientrano nel suo orizzonte di calcolo (\textit{Horizon Effect}).
		\item \textbf{Mancanza di Personalizzazione:} Un avversario che commette sistematicamente errori sulle diagonali viene trattato esattamente come un avversario che commette errori di posizionamento centrale. L'IA non ``impara'' a punire l'errore ricorrente.
		\item \textbf{Over-Defensivismo:} L'IA tende a giocare in modo conservativo, parando minacce che un avversario distratto non avrebbe mai concretizzato, perdendo l'occasione di sviluppare un attacco più aggressivo.
	\end{itemize}
	
	\subsection{Obiettivo del Progetto}
	L'obiettivo primario è la realizzazione di un agente intelligente capace di evolvere la propria strategia in \textit{Real-Time}.
	Il sistema non si limita a calcolare la mossa migliore sulla scacchiera, ma osserva il comportamento dell'utente per costruire un \textbf{Profilo Psicologico} dinamico.
	
	Se il sistema rileva che l'avversario soffre di ``cecità diagonale'' (incapacità di scorgere minacce su linee oblique), l'algoritmo di valutazione euristica deve aumentare il peso specifico dei pattern diagonali, privilegiando le linee di gioco che l'avversario statisticamente non difende.
	Questa architettura sposta il focus dalla pura potenza di calcolo (Deep Blue) alla flessibilità cognitiva, creando un'esperienza di gioco più sfidante e organica.
	
\end{document}